% Chapter 10: The Weak Goldbach Conjecture
\let\LocalMacro\relax

\chapter{The Weak Goldbach Conjecture}

\section{The Problem}

\subsection{Informal}
Pick any odd number bigger than five, like 31 or 999. Can you always write it as the sum of three prime numbers? For example, $31 = 11 + 13 + 7$. Goldbach guessed this was always possible back in 1742, and after nearly three centuries of effort, mathematicians finally proved him right in 2013. The even-number version of the same question, where you only get two primes, remains unsolved to this day.

\begin{historicalbox}
Goldbach's letter to Euler is one of the oldest unsolved problems in mathematics. The strong version remains open after 280 years. The weak version was solved in 2013, just 37 years after the first major breakthrough.
\end{historicalbox}

\subsection{Formal}


\section{Solution}

Who solved it and what techniques were required.

\section{Mathematical Theory}


\subsection{Prerequisites}
This chapter assumes familiarity with:
\begin{itemize}
\item Basic number theory (congruences, primes)
\item Complex analysis (contour integration basics)
\item Fourier analysis (trigonometric sums)
\end{itemize}

\subsection{Vinogradov's Theorem (1937)}

\begin{theorem}[Vinogradov, 1937 \cite{vinogradov1937}]
Every sufficiently large odd integer is the sum of three primes.
\end{theorem}

Vinogradov's proof used the \emph{sieve of Eratosthenes} combined with exponential sums. The key estimate:

\begin{equation}
\sum_{p \leq x} e(\alpha p) \ll x(\log x)^{-A}
\end{equation}

for any $A > 0$, uniformly for $\alpha$ not close to a rational with small denominator. This is the \emph{major arc/minor arc} analysis.

\subsection{The Hardy--Littlewood Circle Method}

The circle method, developed by Hardy and Ramanujan (1918) and refined by Hardy and Littlewood (1923), represents the number of ways to write $n$ as a sum of three primes as an integral:

\begin{equation}
r(n) = \int_0^1 S(\alpha)^3 e(-n\alpha) \, d\alpha
\end{equation}

where $S(\alpha) = \sum_{p \leq n} e(\alpha p)$ is the exponential sum over primes. The interval $[0,1]$ is split into:
\begin{itemize}
\item \textbf{Major arcs}: intervals around rationals $a/q$ with small $q$. Here $S(\alpha)$ is large and can be estimated accurately.
\item \textbf{Minor arcs}: the rest of $[0,1]$. Here $S(\alpha)$ is small, and one needs strong upper bounds.
\end{itemize}

\subsection{The Remaining Challenge}

Vinogradov's theorem only covered \emph{sufficiently large} odd integers. The challenge was to find an explicit bound $N_0$ such that every odd $n > N_0$ is a sum of three primes, and then check all smaller cases computationally.

\section{The Discovery}

\subsection{Ramaré (1995)}

Olivier Ramaré proved that every odd integer $n \geq 1$ is the sum of at most \emph{seven} primes \cite{ramare1995}. This was a significant improvement and reduced the number of primes needed.

\subsection{Helfgott (2013)}

Harald Helfgott completed the proof \cite{helfgott2013}:

\begin{theorem}[Helfgott, 2013 \cite{helfgott2013}]
Every odd integer $n \geq 7$ is the sum of three primes.
\end{theorem}

Helfgott's approach:

\begin{enumerate}
\item \textbf{Improved minor arc bounds}: He developed new estimates for exponential sums on the minor arcs, improving Vinogradov's bounds significantly.
\item \textbf{Explicit major arc analysis}: He computed the major arc contributions with extreme precision.
\item \textbf{Computer verification}: He checked all odd integers up to $10^{30}$ by computer. This was feasible because Ramaré's result reduced the problem to a finite computation.
\end{enumerate}

The proof was verified by a team including David Plante and others using automated theorem proving tools.

\begin{highlightbox}
Helfgott's proof was remarkable because it combined deep analytic number theory with extensive computer verification. The computer-checked over 100 million cases, and the mathematical proof handled the rest.
\end{highlightbox}

\section{Impact}

\begin{itemize}
\item \textbf{Circle method}: Helfgott's improvements to minor arc bounds have applications to other additive problems.
\item \textbf{Computational number theory}: The proof demonstrated the power of combining analytic methods with computation.
\item \textbf{Strong Goldbach}: The weak version's solution brought renewed attention to the strong version, which remains open.
\end{itemize}

\section{Exercises}

\begin{exercise}[Basic]
Verify that every odd integer from 7 to 25 is the sum of three primes. Give explicit decompositions.
\end{exercise}

\begin{exercise}[Intermediate]
Explain why the strong Goldbach conjecture (every even number is the sum of two primes) is harder than the weak version. What makes the two-prime case more difficult?
\end{exercise}

\begin{exercise}[Challenge]
What is the major arc/minor arc decomposition in the circle method? Why do major arcs contribute the main term and minor arcs contribute error terms?
\end{exercise}


\begin{exercise}[Computational]
Write a Python/SageMath script to explore a concept from this chapter. See the appendix for starter code.
\end{exercise}

\section{Timeline}

\begin{tabularx}{\linewidth}{lX}
\toprule
\textbf{Year} & \textbf{Event} \\ \midrule
1742 & Goldbach writes to Euler about his conjecture \\ 
1937 & Vinogradov proves the theorem for sufficiently large $n$ \cite{vinogradov1937} \\ 
1995 & Ramaré proves every odd $n$ is the sum of at most 7 primes \cite{ramare1995} \\ 
2002 & Liu and Wang improve the bound for Vinogradov's theorem \cite{liu2002} \\ 
2013 & Helfgott completes the proof \cite{helfgott2013} \\ \bottomrule
\end{tabularx}

\section{Citations}

\begin{enumerate}
\item Helfgott, H. A. (2013). ``The ternary Goldbach conjecture is true.'' arXiv:1312.7746.
\item Ramaré, O. (1995). ``Every odd number $\geq$ 1 is the sum of at most seven primes.'' Journal de Théorie des Nombres de Bordeaux.
\item Vinogradov, I. M. (1937). ``The method of trigonometrical sums in the theory of prime numbers.'' Doklady Akademii Nauk SSSR.
\item Liu, M. C., \& Wang, T. Z. (2002). ``On the Vinogradov bound in the three primes Goldbach conjecture.'' Acta Arithmetica.
\end{enumerate}
